\section{Complete saturation at the three screened good primes}
\label{sec:137}

We first record the exact passage from component equations to the integral
quotient.  The distinction between the complete equation space and a
control-prime-selected row set is essential in this section.

\begin{lemma}[Good-prime component detection]
\label{lem:good-prime}
Let \(p\nmid|\GammaG|\), suppose \(A_i\otimes\FF_p\) is surjective, and
let \(k/\FF_p\) split \(\GammaG\).  Put
\[
 Q_i=(K_i^{\ZZ}/L_i)\otimes k.
\]
For every absolutely irreducible \(k\GammaG\)-module \(V\), the cokernel
of the \emph{complete} component equation map is dual to
\(\Hom_{k\GammaG}(Q_i,V)\).  Consequently, if every complete component
map has its structural target rank, then
\(\Theta_i\otimes\FF_p=0\), and hence \(\Theta_i[p^\infty]=0\).
\end{lemma}

\begin{proof}
Reduction gives an exact sequence
\[
 0\longrightarrow Q_i
 \longrightarrow k[\cB_i]/(L_i\otimes k)
 \longrightarrow k[\cP]\longrightarrow0.
\]
Maschke's theorem makes \(k\GammaG\) semisimple.  Applying
\(\Hom_{k\GammaG}(-,V)\) identifies the indicated cokernel with the dual
of \(\Hom_{k\GammaG}(Q_i,V)\).  The word ``complete'' is necessary:
replacing all minimum-circuit equations by a row subset selected at a
different prime need not preserve rank after reduction modulo \(p\).
If every absolute type vanishes after a splitting extension, faithful
scalar extension gives \(Q_i=0\).  Finally, finiteness of \(\Theta_i\)
and Nakayama's lemma give \(\Theta_i[p^\infty]=0\).
\end{proof}

At \(p=137\), the hypotheses hold: \(137\nmid204\), every projector
denominator is a unit, and the incidence Smith form has no \(137\)-factor.

\subsection{Why the selected-row obstruction is not a quotient class}

For \(\rh/\cB_1\), a \(1\,312\)-row presentation selected for full rank
at the control prime \(409\) has rank \(1\,311\) at
\[
 \mathfrak p=(137,\zeta_{17}-122).
\]
Its certified residual has rank \(256\) rather than \(257\), and its
local selected-presentation cokernel has length one.  Those calculations
are exact, including the Hensel lift modulo \(137^2\), but they concern
the selected presentation only.

They do not prove a class in \(\Theta_1\): additional minimum-circuit
equations, absent from the \(409\)-selected set, may restore the
structural target rank at \(137\).  This is what happens.

\subsection{The full-circuit restoration}

There are \(11\,274\) \(\GammaG\)-orbits of internal minimum circuits in
\(\cB_1\).  One seed per orbit suffices by equivariance.  At
\(\zeta_{17}=122\), the absolute \(\rh\)-representation has dimension
\(4\); the block multiplicity space has dimension \(1\,317\), while the
point-induced kernel has dimension \(5\).  Hence the maximum circuit
rank is
\[
 1\,317-5=1\,312.
\]

Deterministic sparse elimination over \(\FF_{137}\), in lexicographic
orbit and coordinate order, gives
\[
\begin{array}{c|r}
\text{processed circuit orbits}&\text{rank}\\
\midrule
500&1260\\
1000&1309\\
2000&1311\\
2864&1312.
\end{array}
\]
The restoring source row is coordinate \(0\) of orbit \(2864\)
(zero-based identifier \((2863,0)\)).  An independent reconstruction
forms the selected \(1\,312\times1\,312\) minor and obtains
\[
 \det=56\pmod {137}.
\]
Removing the restoring row leaves rank \(1\,311\), and duplicating a
pivot column makes the determinant zero.

\begin{theorem}[Complete \(\rh/\cB_1\) contribution above \(137\)]
\label{thm:rho137}
The complete minimum-circuit component map for \(\rh/\cB_1\) is full at
every prime above \(137\).  In particular,
\[
 \coker(F_{\rh,1}^{\mathrm{full}})[137^\infty]=0.
\]
The length-one cokernel at \((137,\zeta_{17}-122)\) belongs only to the
control-prime-selected presentation and is not a quotient of
\(\Theta_1\).
\end{theorem}

\begin{proof}
At fifteen of the sixteen split primes, the earlier selected row set
already contains a nonzero maximal minor, so the complete map is full
there.  At the remaining root \(122\), the full-circuit minor above has
nonzero determinant.  Thus every localization above \(137\) is full.
\end{proof}

\subsection{The other five large maps}

For \(\rh\), the sixteen primes above \(137\) are the degree-one roots of
\(\Phi_{17}\).  For \(\rh\om\), put
\[
 k_2=\FF_{137}[y]/(y^2+y+1).
\]
If \(r\) runs over the sixteen order-\(17\) roots in \(\FF_{137}\), then
\[
 \alpha_r=r^6y^2,\qquad
 \alpha_r^3=r,\qquad
 \alpha_r^{17}=y
\]
has order \(51\), and the Frobenius pairs
\(\{\alpha_r,\alpha_r^{137}\}\) represent the sixteen degree-two primes
of \(\ZZ[\zeta_{51}]\) above \(137\).

The exact original-column presentations give:
\[
\begin{array}{c|r|r|c|c}
\text{map}&\text{rows}&\text{columns}&
 \text{residue degrees}&\text{full maximal minors}\\
\midrule
\rh\om/\cB_1&1305&10480&2^{16}&16/16\\
\rh/\cB_2&1314&5276&1^{16}&16/16\\
\rh\om/\cB_2&1304&10472&2^{16}&16/16\\
\rh/\cB_3&1310&5260&1^{16}&16/16\\
\rh\om/\cB_3&1306&10488&2^{16}&16/16.
\end{array}
\]
For the quadratic rows, multiplication by \(a+by\) is expanded as
\[
 \begin{bmatrix}a&-b\\ b&a-b\end{bmatrix}.
\]
A second parser uses the Frobenius-conjugate root, transposes every
matrix, reproduces all \(80\) determinant residues, and kills each
determinant after duplicating one field column.

\begin{theorem}[Complete \(137\)-primary quotients]
\label{thm:theta137}
For the three frozen configurations,
\[
 \Theta_1[137^\infty]
 =\Theta_2[137^\infty]
 =\Theta_3[137^\infty]
 =0.
\]
\end{theorem}

\begin{proof}
\Cref{prop:small-components} excludes \(137\) from all eighteen small
component cokernels.  \Cref{thm:rho137} closes the formerly exceptional
\(\rh/\cB_1\) component, and the eighty nonzero maximal minors close the
other five large maps.  The complete-component criterion in
\cref{lem:good-prime} gives
\(\Theta_i\otimes\FF_{137}=0\) for every \(i\); finiteness then excludes
all \(137\)-primary torsion.
\end{proof}

\begin{remark}[Exact boundary]
\label{rem:137-boundary}
The theorem is complete at the rational prime \(137\).  It does not
classify another rational prime or determine the global Smith invariant
factors of the finite groups \(\Theta_i\).  In particular, singular
minors selected at \(409\) above other candidate primes must first be
audited against the complete circuit equation space.
\end{remark}

\subsection{A second complete prime envelope}

The next selected-presentation signal occurs at \(p=219\,097\).  Since
\[
 219\,097-1=2^3\cdot3^2\cdot17\cdot179,
 \qquad 219\,097\equiv1\pmod {51},
\]
both cyclotomic orders split completely.  Thus a \(\rh\)-map has sixteen
degree-one localizations and a \(\rh\om\)-map has thirty-two.

For \(\rh/\cB_1\), the complete circuit map restores rank \(1\,312\) at
the historical root \(91085\), again at source row \((2863,0)\).  The
independent maximal minor has determinant
\[
 213269\pmod {219097},
\]
and the prefix without the restoring row has rank \(1\,311\).

For the other five maps, exact original-column maximal minors give
\[
\begin{array}{c|r|r|c}
\text{map}&\text{target rank}&\text{localizations}&\text{nonzero minors}\\
\midrule
\rh\om/\cB_1&1305&32&32\\
\rh/\cB_2&1314&16&16\\
\rh\om/\cB_2&1304&32&32\\
\rh/\cB_3&1310&16&16\\
\rh\om/\cB_3&1306&32&32.
\end{array}
\]
An independent parser enumerates the roots from a primitive generator,
reconstructs the transpose of every selected matrix, reproduces all
\(128\) determinants, and obtains determinant zero in all \(128\)
duplicate-column controls.

\begin{theorem}[Complete \(219\,097\)-primary quotients]
\label{thm:theta219097}
For the three frozen configurations,
\[
 \Theta_1[219097^\infty]
 =\Theta_2[219097^\infty]
 =\Theta_3[219097^\infty]
 =0.
\]
\end{theorem}

\begin{proof}
\Cref{prop:small-components} excludes \(219\,097\) from all eighteen small
component cokernels.  The full-circuit determinant closes
\(\rh/\cB_1\), and the \(128\) nonzero maximal minors close the other five
large maps.  The good-prime complete-component criterion
\cref{lem:good-prime} gives \(\Theta_i\otimes\FF_{219097}=0\) for all
three designs; finiteness excludes \(219\,097\)-primary torsion.
\end{proof}

\begin{remark}[Remaining integral boundary]
The theorem removes \(219\,097\) completely from the possible support.  It
does not classify another rational prime or the global Smith invariant
factors.
\end{remark}

\subsection{A third complete prime envelope}

The last selected-presentation signal is
\(p=3\,288\,036\,131\).  Here
\[
 p-1=2\cdot5\cdot17\cdot2243\cdot8623,
 \qquad p\equiv1\pmod {17},\qquad p\equiv2\pmod3.
\]
Thus \(\Phi_{17}\) has sixteen linear factors, while \(\Phi_{51}\)
has sixteen irreducible quadratic factors.

For \(\rh/\cB_1\), the complete circuit map restores rank \(1\,312\) at
the historical root \(1\,120\,906\,769\), again at source row
\((2863,0)\).  The independently reconstructed maximal minor has
\[
 \det=1\,471\,654\,912\pmod {3\,288\,036\,131},
\]
and the prefix without the restoring row has rank \(1\,311\).

For the \(\rh\om\)-maps, use
\[
 k_2=\FF_p[y]/(y^2+y+1),\qquad \alpha_r=r^6y^2,
\]
where \(r\) runs through the order-\(17\) roots.  Then
\(\alpha_r^3=r\), \(\alpha_r^{17}=y\), and the Frobenius pairs give the
sixteen quadratic primes.  The five remaining maps have the exact screen
\[
\begin{array}{c|r|r|c|c}
\text{map}&\text{target rank}&\text{residue degree}&
 \text{localizations}&\text{nonzero minors}\\
\midrule
\rh\om/\cB_1&1305&2&16&16\\
\rh/\cB_2&1314&1&16&16\\
\rh\om/\cB_2&1304&2&16&16\\
\rh/\cB_3&1310&1&16&16\\
\rh\om/\cB_3&1306&2&16&16.
\end{array}
\]
The independent route finds the least primitive generator \(2\), uses
the Frobenius-conjugate quadratic representatives, transposes every
matrix, reproduces all \(80\) determinant residues, and obtains determinant
zero in all \(80\) duplicate-field-column controls.

\begin{theorem}[Complete \(3\,288\,036\,131\)-primary quotients]
\label{thm:theta3288036131}
For the three frozen configurations,
\[
 \Theta_1[3288036131^\infty]
 =\Theta_2[3288036131^\infty]
 =\Theta_3[3288036131^\infty]
 =0.
\]
\end{theorem}

\begin{proof}
\Cref{prop:small-components} excludes \(3\,288\,036\,131\) from all
eighteen small component cokernels.  The full-circuit determinant closes
\(\rh/\cB_1\), and the \(80\) nonzero maximal minors close the other five
large maps.  The complete-component criterion in
\cref{lem:good-prime} gives \(\Theta_i\otimes\FF_p=0\); finiteness excludes
all \(p\)-primary torsion.
\end{proof}

\begin{remark}[Selected signals versus global support]
All three selected good-prime rank losses, at \(137\), \(219\,097\), and
\(3\,288\,036\,131\), disappear for the complete maps.  This finite list
alone does not prove that no untested rational prime divides a global Smith
factor.  The next section supplies that global exclusion away from
\(2,3,17\) for the single complete \(\rh/\cB_1\) component.  The other
five large maps, the support of the complete \(\Theta_i\), and their global
invariant factors remain open.
\end{remark}
